% 四类教学构造。行是二分类任务，列是样本原类别；
% +1 / -1 是该任务的训练标签，点号表示不参与该任务。
\begin{tikzpicture}[x=1cm,y=1cm,font=\small]
  \foreach \shift/\title in {0/一对一（OvO）,5.5/一对其余（OvR）} {
    \begin{scope}[xshift=\shift cm]
      \node[font=\figurefont\small] at (2.4,4.95) {\title};
      \node[font=\footnotesize] at (0.5,4.0) {任务};
      \node[font=\footnotesize] at (2.9,4.5) {样本原类别};
      \foreach \c in {1,2,3,4}
        \node at ({1.65+0.8*(\c-1)},4) {$\c$};
      \draw[FigureGrid] (0,3.75) -- (4.65,3.75);
    \end{scope}
  }
  \foreach \row/\pair/\a/\b/\c/\d in {
    1/{1,2}/+1/-1/\cdot/\cdot,
    2/{1,3}/+1/\cdot/-1/\cdot,
    3/{1,4}/+1/\cdot/\cdot/-1,
    4/{2,3}/\cdot/+1/-1/\cdot,
    5/{2,4}/\cdot/+1/\cdot/-1,
    6/{3,4}/\cdot/\cdot/+1/-1} {
    \pgfmathsetmacro{\rowY}{3.65-0.48*\row}
    \node at (0.5,\rowY) {$h_{\pair}$};
    \node at (1.65,\rowY) {$\a$};
    \node at (2.45,\rowY) {$\b$};
    \node at (3.25,\rowY) {$\c$};
    \node at (4.05,\rowY) {$\d$};
  }
  \begin{scope}[xshift=5.5cm]
    \foreach \row/\a/\b/\c/\d in {
      1/+1/-1/-1/-1,
      2/-1/+1/-1/-1,
      3/-1/-1/+1/-1,
      4/-1/-1/-1/+1} {
      \pgfmathsetmacro{\rowY}{3.65-0.48*\row}
      \node at (0.5,\rowY) {$g_{\row}$};
      \node at (1.65,\rowY) {$\a$};
      \node at (2.45,\rowY) {$\b$};
      \node at (3.25,\rowY) {$\c$};
      \node at (4.05,\rowY) {$\d$};
    }
  \end{scope}
  \draw[book arrow] (2.4,0.35) -- (2.4,-0.15);
  \draw[book arrow] (7.9,0.35) -- (7.9,-0.15);
  \node[book node,align=center,minimum width=4.7cm,font=\small] at (2.4,-0.65)
    {6 个模型各投一票\\选择票数最多的类别};
  \node[book node,align=center,minimum width=4.7cm,font=\small] at (7.9,-0.65)
    {比较 4 个模型的得分\\选择得分最大的类别};
\end{tikzpicture}
